\documentclass{report}
\usepackage{amsmath,amssymb}
\setlength{\parindent}{0mm}
\setlength{\parskip}{1em}
\begin{document}
\begin{center}
\rule{6in}{1pt} \
{\large John Ruge \\
{\bf Hybrid FOSLS/LL* for Navier-Stokes Equations}}

526 UCB \\ University of Colorado \\ Boulder \\ CO 80309-0526
\\
{\tt jruge@colorado.edu}\\
Tom Manteuffel\end{center}

First-Order System Least Squares (FOSLS) is a method for the
discretization of systems of PDEs in which the set of equations is
reformulated as a (generally larger) first-order system. The goal is the
minimization the sum of the squares of the associated residuals (the
FOSLS functional). Properly formulated, this functional is equivalent to
the H1 norm (squared) of the error. This approach has a number of
advantages over traditional methods. In some cases though (e.g. flow
through a long channel), this equivalence can be rather loose, and large
(generally smooth) L2 error can be present even when the functional is
small. In such cases, one approach for linear system is to use FOSLL*,
where the problem is posed to minimize the L2 error over the range of the
adjoint of the operator. This can be problematic for nonlinear problems
though, and a hybrid method was developed. There, both the FOSLS and
FOSLL* functionals are used, along with a connecting term that ties the
two together (by approximately projecting the FOSLL* approximation back
into the original finite element space). This talk covers the approach
taken and results obtained for 2D Stokes and Navier-Stokes flow.
~


\end{document}
